\section{结果与讨论}
\subsection{额外同步记忆带来的误差取舍}
表\ref{tab:accuracy}汇总两个种子的完整结果。
相对\lite，\full 的平均NMAE在Abilene增加0.0382\%，差异很小；
其中第一个种子改善0.2795\%，第二个种子反而恶化0.3601\%，
不能由单个获胜种子推出稳定增量。
在GEANT上，\full 的平均NMAE增加1.4063\%；
两个种子的增加分别为0.6109\%和2.2074\%，
且两个种子下的三个条件都具有相同方向。
因此，在当前主要指标与预算下，新增记忆没有获得默认保留的必要性证据。

平方误差指标呈现另一面。
\full 相对\lite 的平均NRMSE在Abilene降低0.8592\%，在GEANT降低1.2402\%，
四个配对的三条件平均NRMSE均改善。
不过，Abilene第二个种子的“不均匀”条件仍有小幅退化，
因此不能扩大成所有种子、所有条件和所有指标均改善。
现有数据支持增加记忆改变了误差取舍，
而不支持把次要指标的改善替换成主要指标获胜的结论。

SPIN上下文已经聚合整窗可见的时空信息，
Direct又提供了近邻原值的短路径，因而记忆与已有读出可能存在职责重叠。
但这只是解释候选。
Sync读取邻流上下文还可能带来间接邻域信息，且具有不同的矩阵压缩与查询机制。
本轮实验不能区分信息重叠、优化难度和任务中剩余可利用信号不足等原因，
不将当前消融结果解释为已定位某个失败机制。

\subsection{与已有参考的系统比较}
\lite 相对旧SPIN的平均NMAE在Abilene降低1.2235\%，在GEANT降低0.7406\%。
其NRMSE在Abilene降低0.3688\%，在GEANT却增加0.7397\%。
相对Linear，\lite 的平均NMAE分别降低5.0593\%和0.3983\%，
但GEANT上的NRMSE增加1.0591\%。
这说明更简化的组合具备开发NMAE竞争力，同时仍存在平方误差代价。

\full 在Abilene的平均NMAE与NRMSE都低于旧SPIN；
在GEANT，其NMAE高于旧SPIN0.6553\%，NRMSE则低0.5096\%。
旧Direct+Sync在GEANT上保留了更低的平均NRMSE，
Linear也仍是接近主要神经模型的简单参考。
不同参考在绝对误差、平方误差和成本上形成不同取舍。
由于参考模型的深度、优化和检查点选择轨迹不同，
这些比较尚不能把差异全部归因于Direct或某个特定更新规则。

\subsection{推理成本与后续模型选择}
表\ref{tab:cost}显示，\full 相对\lite 的单窗耗时比为4.54和2.36，
批量8时的摊销耗时分别增加64.44\%和62.22\%。
同步记忆的参数增加量固定为24,960，
但参数数量不足以描述两方向扫描及中间张量带来的执行成本。
GEANT的批量8峰值已分配显存增加约27.56\%；
Abilene的对应峰值几乎不变，与共享上下文临时张量占主导的解释一致，
但当前测量没有逐算子定位峰值来源。

\lite 的批量8摊销耗时约为旧SPIN的27.20\%和26.65\%，
由同批量耗时折算，批量推理速度分别约为旧SPIN的3.68和3.75倍。
该优势伴随着一层与四层上下文的结构差异，不能单独归因于Direct。
它也不意味着所有资源都降低：Abilene的参数量比旧SPIN更多，
两个WAN上的峰值显存均仅有极小变化。
既有Direct+Sync在批量8的两个WAN上仍比\lite 更快；
Linear没有参加本轮成本测量，不能据此声称在全部参考中成本最优。

综合预定NMAE、推理耗时和结构复杂度，\lite 是本轮之后更合理的优先简化候选。
若应用更看重平方误差，\full 的结果提供了一个额外成本换取NRMSE改善的实现选择，
但当前证据没有建立这种取舍在独立最终评价上的稳定性。
模型选择应保留任务指标和部署批量，不能只依据参数量或某个获胜数字。
